\documentclass[11pt,a4paper]{article}

\usepackage[T1]{fontenc}
\usepackage[utf8]{inputenc}
\usepackage[a4paper,margin=24mm]{geometry}
\usepackage{amsmath,amssymb}
\usepackage{array}
\usepackage{url}
\Urlmuskip=0mu plus 1mu

\setlength{\parskip}{0.25em}
\setlength{\parindent}{1em}
\emergencystretch=3em
\widowpenalty=10000
\clubpenalty=10000
\displaywidowpenalty=10000

\newcommand{\effect}[1]{\mbox{\texttt{#1}}}
\newcommand{\meet}{\mathbin{\sqcap}}
\newcommand{\seq}{\mathbin{\cdot}}
\newcommand{\sem}[1]{[\![#1]\!]}

\title{Configurable Static Type Checking for Forth:\\
Nine Worked Examples}
\author{Jaanus Pöial\\Tallinn University of Technology}
\date{}

\begin{document}
\maketitle

\begin{abstract}
Forth stack comments document the cells consumed and produced by a word.
Ordinary comments are not checked.  They also do not record how stack words
move or copy individual cells.  This paper presents a configurable source-level
checker.  It represents each phrase by a typed input/output stack contract,
called a stack effect.  A profile supplies the type hierarchy and primitive
word effects.  Symbolic indices record the flow introduced by words such as
\texttt{DUP}, \texttt{SWAP}, and \texttt{OVER}.

The analysis uses three partial operations.  Composition handles word
sequences.  Branches keep only facts that hold on both paths.  A finite
one-pass-versus-two-pass comparison handles repetition.  Nine small Forth
definitions illustrate these operations.  They cover type refinement, cell
identity, hidden stack prefixes, branch merging, and local loop stability.
The checker reproduces all nine reported effects without executing concrete
values.  Its conclusions depend on the supplied primitive specifications.
The branch summary also forgets path-specific facts.  The loop summary does
not cover every possible iteration count.
\end{abstract}

\noindent\textbf{Keywords:} Forth; static type checking; stack effects;
symbolic execution; subtyping; concatenative languages.

\section{Introduction}

A Forth word is a named operation.  Here, a cell is one single-cell stack
item.  The data stack is the ordered sequence of those items.  A word normally
receives operands from that stack and returns its results to the same stack.
The conventional notation
\begin{verbatim}
+   ( n n -- n )
\end{verbatim}
places the top of stack on the right and says that \texttt{+} consumes two
numbers and produces one.  Such comments are valuable interfaces, but a Forth
system usually consumes them as source text rather than treating them as
proof obligations.  A comment can therefore be missing, stale, or semantically
stronger than the implementation.

Consider the familiar definition
\begin{verbatim}
: SQUARED  ( n -- n^2 )  DUP * ;
\end{verbatim}
The checker can derive \effect{( n -- n )} from the declarations of
\texttt{DUP} and \texttt{*}.  It can verify the number and identity of the
cells used by the phrase.  It does not prove the mathematical assertion that
the result is the square of the input.  This distinction---typed stack
consistency rather than arbitrary functional correctness---sets the scope of
the method.

\begin{samepage}
The checker reads three text inputs:
\begin{enumerate}
  \item a file that defines types, aliases, and ordering relations;
  \item a file that defines word effects and source-reading rules;
  \item a Forth source file whose colon definitions are to be inferred.
\end{enumerate}
A \emph{primitive word} is a word whose effect comes directly from the second
file.  A \emph{colon definition} begins with \texttt{:} and ends with
\texttt{;}.  The word that it defines receives an inferred effect from the
source between those delimiters.
\end{samepage}
The checker has Java and Gforth implementations.  Both consume the same three
files.  An analysis dictionary maps each known word to one stack effect.  At
\texttt{;}, the inferred body effect is added to this dictionary.  The checker
does not execute the Forth definition.

The design builds on earlier work on stack effects and symbolic type variables
\cite{poial1994,poial1992,poial2002,poial2006,poial2008}.  This paper gives a
self-contained account of the operations used by the nine examples.  It
explains four points:
\begin{itemize}
  \item how a later word can constrain an earlier input;
  \item how two branch effects are combined;
  \item why an effect displays only the top cells that it touches;
  \item why the loop check compares one pass with two passes.
\end{itemize}
Sections~\ref{sec:model}--\ref{sec:control} define the complete analysis used
here.  Sections~\ref{sec:sequence-examples}--\ref{sec:loop-examples} derive all
nine results.  Appendices A--C reproduce the source and every profile entry
needed for those derivations.  No other project document is required.

\section{Typed Symbolic Stack Effects}
\label{sec:model}

\subsection{Types and compatibility}

Let $T$ be a finite set of abstract cell types with a precision order
$\preceq$.  The relation $t\preceq u$ means that $t$ is more specific than
$u$.  A value of type $t$ may therefore be used where $u$ is required.  This
order is the profile's subtype relation.  The Forth-2012-like profile used for the examples contains, among
other edges, the transitive chains
\[
 \texttt{a-addr}\preceq\texttt{c-addr}\preceq
 \texttt{addr}\preceq\texttt{u}\preceq\texttt{n}\preceq\texttt{x},
\]
\[
 \texttt{char}\preceq\texttt{+n}\preceq
 \texttt{u}\preceq\texttt{n}\preceq\texttt{x},
 \qquad \texttt{flag}\preceq\texttt{x}.
\]
The flag type is deliberately separate from the numeric chain in this profile.
Another profile may classify flags as numbers.  The hierarchy is therefore a
policy for one target vocabulary.  It is not a universal claim about Forth
cells.  Names such as \texttt{x}, \texttt{X}, and \texttt{cell} may be aliases
for one type.

Two types are \emph{comparable} when one is below the other in $\preceq$.
Comparable types are compatible.  Matching them retains the more specific
type:
\[
 \operatorname{match}(t,u)=
 \begin{cases}
 t,&t\preceq u,\\
 u,&u\preceq t,\\
 \text{undefined},&\text{otherwise}.
 \end{cases}
\]
The implementation does not search for a third common subtype when $t$ and
$u$ are incomparable.  An undefined match is a type clash.

\subsection{Symbols, effects, and the implicit frame}

A symbolic cell is a type with an optional positive identity index, written
$t[i]$.  Equal indices denote the same abstract origin.  Different indices do
not claim that the concrete values differ.  Every unindexed occurrence starts
as an independent symbol.  An unindexed output is therefore a fresh abstract
result.  Indices exist only in the checker.  They are not addresses or runtime
tags.

\begin{samepage}
The declarations of three stack words illustrate the notation:
\begin{verbatim}
DUP   ( x[1] -- x[1] x[1] )
SWAP  ( x[2] x[1] -- x[1] x[2] )
OVER  ( x[2] x[1] -- x[2] x[1] x[2] )
\end{verbatim}
Before each word occurrence is analyzed, its local indices receive fresh
numbers.  Two primitive calls therefore cannot share an index by accident.
After analysis, \emph{normalization} renumbers the surviving indices from one
upward.  Renumbering changes only their printed names, not their relationships.
\end{samepage}

An effect $e=(L--R)$ contains two finite symbolic stacks, with their tops on
the right.  A stack \emph{suffix} is the group of top cells displayed by the
effect.  Deeper cells form an implicit \emph{frame}.  For every frame $F$, the
contract acts as
\[
                  F\,L \longmapsto F\,R.
\]
For example, \effect{( a-addr -- x )} for \texttt{@} does not require the
entire data stack to contain one cell.  It requires only that the top cell be
an aligned address; any deeper $F$ passes through unchanged.  The identity
effect is \effect{( -- )}.  Write $\epsilon$ for the empty stack.  A clash is
a separate failure result and must not be confused with the identity effect.
The set of symbolic effects is the checker's \emph{abstract domain}.  An
effect operation is \emph{partial} when some arguments produce no effect.  In
those cases, this checker reports a clash.

\section{The Three Effect Operations}
\label{sec:operations}

\subsection{Sequential composition}

Let $p=(L_p--R_p)$ and $q=(L_q--R_q)$.  The expression $p\seq q$ means: apply
$p$, then apply $q$.  Composition compares the top output of $p$ with the top
input of $q$.  It then works inward.
For every compatible pair, the checker:
\begin{enumerate}
  \item selects the more specific boundary type;
  \item creates a fresh shared symbol of that type;
  \item substitutes it for both old symbols throughout the working effects;
  \item removes the matched output/input pair and continues.
\end{enumerate}
If $p$ runs out of outputs first, $q$ still needs cells from the implicit
frame.  Those cells become inputs of the composite.  If $q$ runs out of inputs
first, unused outputs of $p$ remain below the outputs of $q$.  This is exactly
the effect of applying the two suffix transformations in order.

Substitution through the whole effect is essential.  A boundary constraint
must reach every copy with the same index.  A later \texttt{@} can therefore
refine a cell that passed through \texttt{SWAP} and \texttt{DUP}.
Incompatible types make composition undefined.  A word sequence is reduced by
repeated composition from left to right.  The empty sequence has the identity
effect.  This operation extends earlier untyped stack-effect composition
\cite{poial1994}.

\subsection{Common branch contracts}
\label{sec:branch-op}

Alternative paths use a second partial operation, written $p\meet q$.  Call it
the \emph{common contract}.  It contains only facts justified for both paths.
It does not store a separate effect for each path.  Positions are first aligned
from the top of each stack.

At aligned inputs, the operation keeps the more specific comparable type.  The
input to the complete conditional must satisfy both paths.  At aligned outputs,
it keeps the more general comparable type.  Later code may use only a type
guaranteed by both paths.  This opposite treatment of inputs and outputs is
called \emph{variance}.  An output index survives only if both paths preserve
the same input-to-output relationship.  A branch-dependent permutation is
therefore forgotten.

Displayed effects may have unequal depths because both omit an unchanged frame.
Suppose $p$ shows $k$ more input cells than $q$.  A merge is considered only if
$p$ also shows exactly $k$ more output cells.  The missing $k$-cell prefix of
$q$ is treated as unchanged, after which active suffixes are aligned.  Formally,
a necessary shape condition is
\[
 |L_p|-|L_q| = |R_p|-|R_q|.
\]
A different net stack change or an incomparable aligned type rejects the join.

Write $\sem{P}$ for the inferred effect of a source segment $P$.  For an
\texttt{IF}, the declared selector effect is composed separately:
\[
 \sem{\texttt{IF }P\texttt{ ELSE }Q\texttt{ FI}}
   = (\texttt{flag}--)\seq
     (\sem{P}\meet\sem{Q}).
\]
When \texttt{ELSE} is absent, $Q$ is the identity effect.

\subsection{Finite repetition}
\label{sec:repeat-op}

For a repeated phrase with one-pass effect $p$, the prototype computes
\[
                  \pi_2(p)=p\meet(p\seq p),
\]
provided both the composition and common-contract operations are defined.
This one-versus-two comparison is finite and terminating.  It detects many
bodies that grow, shrink, or change the stack incompatibly on a second pass.

The notation has a narrow meaning.  It does not summarize every possible
iteration count.  One and two passes can agree after information has been
lost.  A later pass may then remove another correlation.  A full loop analysis
would propagate the loop-header state until that state stops changing.
$\pi_2$ does not perform that process.  It also proves neither termination nor
properties of concrete values.  It is only a local stack-stability test.

\section{Configuration and Experimental Method}
\label{sec:control}

The run uses three files: a type profile, a word profile, and the example
source.  Appendices A--C reproduce every line from those files that affects the
nine definitions.  Table~\ref{tab:primitives} summarizes the primitive
contracts.  The integer-literal entry classifies tokens such as \texttt{0},
\texttt{1}, and \texttt{5}.  It does not evaluate their values.

\clearpage
\begin{table}[h]
\centering
\small
\begin{tabular}{@{}lll@{}}
\hline
Class & Word & Declared effect \\
\hline
Literals & integer, \texttt{FALSE} & \texttt{( -- n )}, \texttt{( -- flag )} \\
Stack & \texttt{DUP} & \texttt{( x[1] -- x[1] x[1] )} \\
      & \texttt{SWAP} & \texttt{( x[2] x[1] -- x[1] x[2] )} \\
      & \texttt{OVER} & \texttt{( x[2] x[1] -- x[2] x[1] x[2] )} \\
      & \texttt{ROT} & \texttt{( x[3] x[2] x[1] -- x[2] x[1] x[3] )} \\
Numeric & \texttt{+}, \texttt{-} & \texttt{( n n -- n )} \\
Logic & \texttt{OR} & \texttt{( flag flag -- flag )} \\
      & \texttt{INVERT} & \texttt{( flag -- flag )} \\
Memory & \texttt{@} & \texttt{( a-addr -- x )} \\
       & \texttt{C@} & \texttt{( c-addr -- char )} \\
Input & \texttt{KEY} & \texttt{( -- char )} \\
Control & \texttt{IF}, \texttt{WHILE}, \texttt{UNTIL}
        & \texttt{( flag -- )} \\
\hline
\end{tabular}
\caption{Primitive effects used by the examples.}
\label{tab:primitives}
\end{table}

The presentation shows \texttt{0=} as \effect{( n -- flag )}.  The word
profile in Appendix C uses the broader effect \effect{( x -- flag )}.  This
difference does not change a reported result.  In \texttt{test8}, \texttt{KEY}
produces \texttt{char}.  The relation
$\texttt{char}\preceq\texttt{n}\preceq\texttt{x}$ makes both declarations
applicable.

The word profile also describes control syntax.  Angle brackets name captured
source segments.  Square brackets mark an optional segment.  Effect-recipe
lines are composed from top to bottom.  A bare control word contributes its
declared effect.  \texttt{either} applies the common-contract operation
$\meet$.  \texttt{repeat} applies the finite operation $\pi_2$.

The conditional declaration is:
\begin{verbatim}
syntax: IF <then branch> [ELSE <else branch>] FI
  effect:
    IF
    either <then branch> <else branch>
\end{verbatim}
\medskip
\noindent For loops, the relevant recipes are:
\begin{verbatim}
syntax: BEGIN <loop prefix> WHILE <loop body> REPEAT
  effect:
    repeat <loop prefix> WHILE
    repeat <loop body>

syntax: BEGIN <loop body> UNTIL
  effect:
    repeat <loop body> UNTIL
\end{verbatim}
The syntax declarations split source into segments.  The effect recipes then
combine those segments with the three operations already defined.

The exact definitions were checked by loading the three files.  The resulting
word effects were read from the definition log.  One reproducible invocation
is:
\begin{verbatim}
./run-evaluator-gforth.sh --prog examples.txt
\end{verbatim}
Equivalently, after compiling the Java sources:
\begin{verbatim}
javac evaluator/*.java
java evaluator.Evaluator \
  --types forth2012types.txt \
  --specs forth2012specs.txt \
  --prog examples.txt
\end{verbatim}
\clearpage
All nine definitions are accepted.  Table~\ref{tab:results} reproduces their
inferred contracts.  The next sections derive each result.

\begin{table}[h]
\centering
\small
\begin{tabular}{@{}lll@{}}
\hline
Word & Inferred effect & Principal mechanism \\
\hline
\texttt{test1} & \texttt{( -- n )} & numeric sequence \\
\texttt{test2} & \texttt{( flag flag -- flag flag )} & flag propagation \\
\texttt{test3} & \texttt{( a-addr[2] x[1] -- x[1] a-addr[2] x )}
                & identity and backward refinement \\
\texttt{test4} & \texttt{( n flag -- n )} & implicit identity branch \\
\texttt{test5} & \texttt{( x[1] flag -- x[1] x )} & output widening \\
\texttt{test6} & \texttt{( x x a-addr flag -- x x x )}
                & path-dependent permutation loss \\
\texttt{test7} & \texttt{( flag flag -- flag flag )} & two-part repetition \\
\texttt{test8} & \texttt{( -- )} & balanced test loop \\
\texttt{test9} & \texttt{( a-addr flag -- x )} & input/output variance \\
\hline
\end{tabular}
\caption{Exact inferred effects recorded in \texttt{examples.txt.log}.}
\label{tab:results}
\end{table}

\section{Straight-Line Examples}
\label{sec:sequence-examples}

\subsection{\texttt{test1}: stack shape without value evaluation}

\begin{quote}
\texttt{: test1  3 4 5 - + ;}
\end{quote}
Every literal has effect \effect{( -- n )}.  Naming fresh abstract results
$a$, $b$, and $c$, the stack evolves as
\[
 \epsilon
 \longrightarrow n_a
 \longrightarrow n_a\;n_b
 \longrightarrow n_a\;n_b\;n_c
 \longrightarrow n_a\;n_d
 \longrightarrow n_e.
\]
Subtraction consumes $n_b,n_c$ and produces a fresh $n_d$.  Addition consumes
$n_a,n_d$ and produces $n_e$.  No input is required.  One number is left:
\[
                      \texttt{test1}\;( -- n ).
\]
Concrete execution would calculate $4-5=-1$ and then $3+(-1)=2$.  The checker
does not perform that calculation.  Its domain contains the type \texttt{n},
but no integer constants or arithmetic expressions.  It establishes only
\effect{( -- n )}.

\subsection{\texttt{test2}: a generic shuffle specialized to flags}

\begin{quote}
\texttt{: test2  OR FALSE SWAP ;}
\end{quote}
\texttt{OR} first requires two flags and leaves one fresh flag.  \texttt{FALSE}
adds another flag.  The generic inputs of \texttt{SWAP} then match these two
results.  Both inputs are refined from \texttt{x} to \texttt{flag}:
\[
 \texttt{flag}_a\;\texttt{flag}_b
 \xrightarrow{\texttt{OR}}
 \texttt{flag}_c
 \xrightarrow{\texttt{FALSE}}
 \texttt{flag}_c\;\texttt{flag}_d
 \xrightarrow{\texttt{SWAP}}
 \texttt{flag}_d\;\texttt{flag}_c.
\]
The inputs are consumed by \texttt{OR}; neither output is correlated with an
original input.  The reported interface therefore needs no visible indices:
\[
           \texttt{test2}\;(\texttt{flag flag -- flag flag}).
\]
Although one output is produced by the word named \texttt{FALSE}, the checker
has no constant-propagation domain and promises only its declared type.

\subsection{\texttt{test3}: backward refinement through \texttt{SWAP} and
\texttt{DUP}}

\begin{quote}
\texttt{: test3  SWAP DUP @ ;}
\end{quote}
Let the initial suffix be $A\;B$, with $B$ topmost.  \texttt{SWAP} produces
$B\;A$ and \texttt{DUP} produces $B\;A\;A$.  The final \texttt{@} requires
the top copy of $A$ to be an \texttt{a-addr}.  Because both copies carry the
same symbolic index, the substitution also refines the untouched copy and the
original input:
\[
 A\;B
 \xrightarrow{\texttt{SWAP}} B\;A
 \xrightarrow{\texttt{DUP}} B\;A\;A
 \xrightarrow{\texttt{@}} B\;A\;X_{\mathrm{fresh}}.
\]
After normalized numbering, $A$ is \texttt{a-addr[2]} and $B$ is
\texttt{x[1]}.  Hence
\[
\begin{split}
 \texttt{test3}\;(&\texttt{a-addr[2] x[1]}\\
                  &\texttt{-- x[1] a-addr[2] x}).
\end{split}
\]
This is the key identity-flow example.  Tracking only stack depth would miss
the address requirement.  Tracking types without indices would also lose
information.  Such an analysis could refine only the consumed copy.  It could
not identify the surviving copy as the same aligned address.

\section{Branch Examples}
\label{sec:branch-examples}

\subsection{\texttt{test4}: an explicit effect versus an implicit frame}

\begin{quote}
\texttt{: test4  IF 1 + FI ;}
\end{quote}
The true branch \texttt{1 +} pushes an \texttt{n}, then needs one additional
\texttt{n} below it for \texttt{+}.  Its inferred effect is
\effect{( n -- n )}.  The absent \texttt{ELSE} branch has the identity effect
\effect{( -- )}.  These displayed depths differ, but the empty branch preserves
an arbitrary deeper frame.  Making one frame cell explicit gives the
alignment
\[
\begin{array}{rcl}
 \text{true}  &:& (\texttt{n -- n}),\\
 \text{false} &:& (\underline{\texttt{n}} --
                       \underline{\texttt{n}}).
\end{array}
\]
Both paths preserve a one-cell numeric stack shape.  They do not preserve the
same value.  The false path returns the original number.  The true path returns
the result of addition.  No output identity is claimed.
After the leading \texttt{IF} consumes its selector, the complete definition is
\[
                    \texttt{test4}\;(\texttt{n flag -- n}).
\]
This example explains why comparing only the explicitly printed sides would be
too strict.  Stack effects denote transformations under an implicit frame.

\subsection{\texttt{test5}: a common prefix survives, a branch-dependent
identity does not}

\begin{quote}
\texttt{: test5  IF DUP ELSE 0 FI ;}
\end{quote}
The true branch is
\effect{( x[1] -- x[1] x[1] )}.  The false branch explicitly has
\effect{( -- n )}, but it leaves the deeper frame unchanged.  Exposing one
frame cell aligns the branches as follows:
\[
\begin{array}{rcl}
 \text{true}  &:& (\texttt{x[1] -- x[1] x[1]}),\\
 \text{false} &:& (\texttt{x[1] -- x[1] n}).
\end{array}
\]
The lower output is the original input on both paths, so its index survives.
The top output is the same input on the true path but a fresh numeric literal
on the false path.  Its path-independent type is the more general type
\texttt{x}.  It cannot retain the input index.  The branch contract is thus
\[
                       (\texttt{x[1] -- x[1] x}).
\]
Composing the selector yields
\[
             \texttt{test5}\;(\texttt{x[1] flag -- x[1] x}).
\]
The result loses two facts.  First, \texttt{n} is widened to \texttt{x} at the
differing output.  Second, a correlation valid on only one path is discarded.

\subsection{\texttt{test6}: compatible shape with no common permutation}

\begin{quote}
\texttt{: test6  IF ROT ELSE @ FI ;}
\end{quote}
Write the pre-branch suffix as $A\;B\;C$, with $C$ topmost.  The true branch
rotates it:
\[
                         A\;B\;C \longmapsto B\;C\;A.
\]
The false branch \texttt{@} displays only one input and one output.  Its
implicit two-cell frame makes the comparable transformation
\[
                         A\;B\;C \longmapsto A\;B\;D,
\]
where $C$ must be \texttt{a-addr} and $D$ is a fresh \texttt{x}.  At the input,
\texttt{ROT}'s generic requirement and \texttt{@}'s address requirement meet
at the more specific \texttt{a-addr}.  The input contract is consequently
\effect{( x x a-addr -- ... )}.

The output positions have the path pairs $(B,A)$, $(C,B)$, and $(A,D)$.  None
of these establishes one common source identity.  Their only common output
type is \texttt{x}, so the branch contract is
\[
                  (\texttt{x x a-addr -- x x x}).
\]
After selector consumption, the inferred definition is
\[
       \texttt{test6}\;(\texttt{x x a-addr flag -- x x x}).
\]
The checker accepts the equal net stack change.  It forgets which input reaches
which output.  An analysis that stores each path separately could retain both
permutations.  This checker stores one common contract.

\subsection{\texttt{test9}: input and output variance in isolation}

\begin{quote}
\texttt{: test9  IF @ ELSE C@ FI ;}
\end{quote}
The two branch effects are
\[
 (\texttt{a-addr -- x})
 \qquad\text{and}\qquad
 (\texttt{c-addr -- char}).
\]
An input to the whole conditional must satisfy both address requirements.
Since \texttt{a-addr}$\preceq$\texttt{c-addr}, the common input is the more
specific \texttt{a-addr}.  An output client, however, needs a guarantee valid
for either \texttt{x} or \texttt{char}.  Since
\texttt{char}$\preceq$\texttt{x}, that guarantee is the more general
\texttt{x}.  Thus
\[
 (\texttt{a-addr -- x})\meet(\texttt{c-addr -- char})
       =(\texttt{a-addr -- x}),
\]
and selector composition gives
\[
               \texttt{test9}\;(\texttt{a-addr flag -- x}).
\]
This small example is the clearest demonstration of contract variance:
branch inputs are narrowed, whereas outputs are widened.

\section{Repetition Examples}
\label{sec:loop-examples}

\subsection{\texttt{test7}: one-versus-two exposes both flag requirements}

\begin{quote}
\texttt{: test7  BEGIN SWAP OVER WHILE INVERT REPEAT ;}
\end{quote}
The profile divides this structure into the repeated prefix
\texttt{SWAP OVER WHILE} and the repeated body \texttt{INVERT}.  Consider one
prefix pass on $A\;B$:
\[
 A\;B
 \xrightarrow{\texttt{SWAP}} B\;A
 \xrightarrow{\texttt{OVER}} B\;A\;B
 \xrightarrow{\texttt{WHILE}} B\;A.
\]
\texttt{WHILE} consumes the duplicate of the original top cell, so $B$ must be
a flag.  A second prefix pass transforms $B\;A$ back to $A\;B$ and tests $A$.
Consequently the two-pass effect requires both original cells to be flags.
The common summary preserves two flag-typed positions.  It cannot promise one
permutation.  One pass swaps the cells, while two passes restore them.  The
summary is therefore
\[
                        (\texttt{flag flag -- flag flag}).
\]
The second repeated region, \texttt{INVERT}, has the stable type effect
\effect{( flag -- flag )}.  Composing it with the prefix summary leaves the
same two-cell interface, giving
\[
          \texttt{test7}\;(\texttt{flag flag -- flag flag}).
\]
This is a useful rejection test for changing stack shape and incompatible
second-pass types.  It is not a proof that the reported identity-free summary
characterizes every number of loop iterations or that the loop terminates.

\subsection{\texttt{test8}: a data-stack-neutral loop pass}

\begin{quote}
\texttt{: test8  BEGIN KEY 0= UNTIL ;}
\end{quote}
A complete repeated pass has the data-stack trace
\[
 \epsilon
 \xrightarrow{\texttt{KEY}} \texttt{char}
 \xrightarrow{\texttt{0=}} \texttt{flag}
 \xrightarrow{\texttt{UNTIL}} \epsilon.
\]
The displayed numeric contract for \texttt{0=} accepts \texttt{char} because
\texttt{char}$\preceq$\texttt{n}.  The broader \texttt{x} contract also
accepts it.  \texttt{UNTIL} consumes the resulting flag.  One pass has the
identity data-stack effect.  Two passes have the same effect.  Their common
contract is
\[
                         \texttt{test8}\;(\texttt{--}).
\]
This conclusion concerns only the modeled data stack.  \texttt{KEY} still has
an I/O side effect.  The result does not promise a key that meets the exit
condition.  It also does not prove termination.

\section{Discussion}

\subsection{What the examples establish}

The nine examples isolate complementary properties of the symbolic domain.
\texttt{test1} shows that the checker infers shape and cell type without
constant evaluation.  \texttt{test2} shows subtype specialization of a generic
stack word.  \texttt{test3} demonstrates bidirectional information flow:
execution order is left to right, but a late requirement propagates backward
through shared symbolic identities.

The branch examples show why an effect is more than a pair of stack heights.
\texttt{test4} depends on the implicit frame.  \texttt{test5} retains only an
identity common to both paths.  \texttt{test6} loses a branch-dependent
permutation.  \texttt{test9} narrows its input and widens its output.  The loop
examples show two uses of the finite rule.  A second pass refines another cell
in \texttt{test7}.  A balanced pass becomes identity in \texttt{test8}.

All reported outputs are inferred from primitive contracts.  The illustrative
comments in the presentation are not declarations.  This rule prevents a stale
comment from changing inference.  Semantic labels such as \texttt{value} or
``address contents'' are therefore outside the domain.  A richer profile would
be needed to represent them.

\subsection{Configuration as an analysis parameter}

The same source may receive a different verdict under another valid profile.
For example, one profile may classify \texttt{char} as numeric.  It accepts
the displayed \texttt{KEY 0=}.  Another profile may keep \texttt{char} and
\texttt{n} unrelated.  A numeric-only \texttt{0=} then rejects the phrase.
The target Forth system determines which policy is appropriate.  The profile
makes that policy reviewable.

External control declarations provide similar flexibility.  They also enlarge
the trusted base.  A recipe might omit a selector effect.  It might also put
\texttt{UNTIL} outside its repeated segment.  Either error can produce a wrong
summary even when the three operations are correct.  The configuration should
therefore be reviewed with the checked code.

\section{Validity and Limitations}
\label{sec:limits}

Acceptance has a precise but bounded interpretation:
\begin{itemize}
  \item every checked linear boundary has enough modeled cells of comparable
        types;
  \item every checked branch admits one common contract under the implemented
        alignment and variance rules;
  \item every checked repetition admits the finite $p\meet(p\seq p)$ summary;
  \item parser roles, delimiters, and compilation state fit the active profile.
\end{itemize}
It does not imply full Forth program correctness.  The principal threats to a
stronger interpretation are as follows.

\paragraph{Trusted declarations.}
Primitive and parser specifications are assumptions.  If \texttt{@},
\texttt{KEY}, or a control word is declared incorrectly, analysis of its
clients cannot make the result reliable.

\paragraph{Single effects and comparable types.}
The dictionary stores one effect per primitive.  Branch merging also returns
one common contract.  It cannot retain a separate permutation for each path.
Aligned types must be directly comparable.  The checker does not search for a
third common type.

\paragraph{Correlation precision.}
Normalization and single-contract merging can discard symbolic identities.
This loss is conservative for the displayed examples but can make later
analysis less precise.  Internal correlations should ideally be retained and
compacted only for presentation.

\paragraph{Finite loop approximation.}
The one-versus-two comparison does not prove a result for every iteration
count.  A stronger analysis would maintain an abstract stack at the loop
header.  It would propagate the body until that stack stopped changing.
Another design would require a declared loop interface and check every back
edge against it \cite{cousot1977}.

\paragraph{Forth language coverage.}
Return-stack manipulation, unrestricted \texttt{EXECUTE}, recursion,
non-local \texttt{THROW}, native defining words, self-modifying dictionaries,
and unrestricted metaprogramming require additional abstractions.  The checker
should complement runtime tests, code review, and target-system knowledge.

The examples are explanatory cases.  They are not a performance benchmark.
They also do not prove that every accepted program is safe.  Each case makes
one operation visible.  The exact nine contracts can be compared across the
Java and Gforth implementations.

\section{Related Work}

Forth has long used stack-effect notation as an interface convention.  The
Forth standard defines data types and control semantics \cite{forth2012}.  It
does not make each source comment a static judgment.  Earlier work developed
algebraic composition and multiple effects for control flow
\cite{poial1994,poial1992}.  Later work added typed symbolic variables and
conservative control-flow operations \cite{poial2002,poial2006}.  A Java
framework made these ideas executable \cite{poial2008}.  The checker described
here also loads parser and control behavior from a profile.

Stoddart and Knaggs also studied type inference for stack languages
\cite{stoddart1991}.  Saabas and Uustalu gave compositional type systems for
stack-based low-level code \cite{saabas2006}.  JVM bytecode verification also
propagates stack types through control-flow joins \cite{leroy2003}.  Those
settings have more fixed instruction and control languages.  Forth source can
contain parser words and dictionary extensions.  The profile handles that
variability, but it becomes part of the trusted input.

Abstract interpretation offers a stronger treatment of loops
\cite{cousot1977}.  It propagates abstract states until they stop changing.
The one-versus-two rule performs only the first bounded comparison.  An
inferred loop effect must be read with that distinction in mind.  Every
bibliography entry below includes a public full-text link.  The links were
checked on 30 August 2026.

\section{Conclusion}
\leavevmode\par\vspace{-\baselineskip}
Typed symbolic stack effects turn a familiar Forth notation into executable,
configurable documentation.  Sequential composition checks touching stack
boundaries and propagates refined types through shared identities.  Branch
merging computes one conservative contract, narrowing inputs, widening outputs,
and preserving only path-independent correlations.  Repetition compares one
pass with two to obtain a useful but explicitly finite stability summary.

The nine examples make these mechanisms concrete.  Three linear definitions
show type and identity propagation.  Four conditionals show implicit frames,
input narrowing, output widening, and information loss.  Two loops show what
the bounded repetition rule accepts.  The profile reproduces every inferred
effect.  The result is a practical checker for the configured subset of Forth.
It is not a proof of functional behavior or termination.  Stronger assurance
requires validated primitives, separate path effects, preserved identity
classes, and full loop analysis.

\appendix
\section{Exact Example Source and Inferred Log}

\begin{samepage}
The source file used for Table~\ref{tab:results} is:
\begin{verbatim}
: test1 3 4 5 - + ;
: test2 OR FALSE SWAP ;
: test3 SWAP DUP @ ;
: test4 IF 1 + FI ;
: test5 IF DUP ELSE 0 FI ;
: test6 IF ROT ELSE @ FI ;
: test7 BEGIN SWAP OVER WHILE INVERT REPEAT ;
: test8 BEGIN KEY 0= UNTIL ;
: test9 IF @ ELSE C@ FI ;
\end{verbatim}
\end{samepage}
\begin{samepage}
The generated definition log is:
\begin{verbatim}
test1 ( -- n )
test2 ( flag flag -- flag flag )
test3 ( a-addr[2] x[1] -- x[1] a-addr[2] x )
test4 ( n flag -- n )
test5 ( x[1] flag -- x[1] x )
test6 ( x x a-addr flag -- x x x )
test7 ( flag flag -- flag flag )
test8 ( -- )
test9 ( a-addr flag -- x )
\end{verbatim}
Whitespace in the machine log is presentation-only; the symbolic contracts are
as shown.
\end{samepage}

\section{Relevant Type-Profile Declarations}
\label{app:types}

The excerpt below contains the types, aliases, and subtype edges needed by the
nine examples.  It is copied from \path{forth2012types.txt}.  Unused profile
types, such as execution tokens, file identifiers, and double-cell numbers,
are omitted.
\begin{verbatim}
type x X cell
type n number signed
type +n nonnegative
type u unsigned
type char c
type flag boolean
type addr
type c-addr ca
type a-addr aa

rel +n < n
rel +n < u
rel u < n
rel n < x
rel u < x
rel char < +n
rel flag < x
rel addr < u
rel c-addr < addr
rel a-addr < c-addr
\end{verbatim}
The loader adds every relation implied by a chain.  For example, it derives
\texttt{char}$\preceq$\texttt{n}$\preceq$\texttt{x}.  It also derives
\texttt{a-addr}$\preceq$\texttt{c-addr}$\preceq$\texttt{x}.
The type \texttt{flag} remains outside the numeric chain.

\section{Relevant Word and Control Specifications}
\label{app:specs}

This appendix selects the entries from \texttt{forth2012specs.txt} used to
parse and infer \texttt{examples.txt}.  The text of each declaration is
unchanged; unrelated vocabulary entries are omitted.

\subsection{Stack words, definition syntax, and literals}
\begin{verbatim}
OVER ( x[2] x[1] -- x[2] x[1] x[2] )
SWAP ( x[2] x[1] -- x[1] x[2] )
DUP ( x[1] -- x[1] x[1] )
ROT ( x[3] x[2] x[1] -- x[2] x[1] x[3] )

: parse word define colon ( -- )
; control end ( -- )
literal integer ( -- n )
FALSE ( -- flag )
\end{verbatim}
The first four declarations expose symbolic identity.  The source-reading
markers on \texttt{:} and \texttt{;} start and close a colon definition.  The
literal class supplies the effect of tokens such as \texttt{0}, \texttt{1},
and \texttt{5}.

\subsection{Conditional control structure}
\begin{verbatim}
IF control if ( flag -- )
ELSE control else ( -- )
FI control fi ( -- )

syntax: IF <then branch> [ELSE <else branch>] FI
  effect:
    IF
    either <then branch> <else branch>
\end{verbatim}
The control roles delimit the captured source regions.  The effect recipe first
accounts for selector consumption and then applies the common-contract
operation to the two branches.

\subsection{Loop control structures}
\begin{verbatim}
BEGIN control begin ( -- )
WHILE control while ( flag -- )
REPEAT control repeat ( -- )

syntax: BEGIN <loop prefix> WHILE <loop body> REPEAT
  effect:
    repeat <loop prefix> WHILE
    repeat <loop body>
\end{verbatim}
\begin{samepage}
\begin{verbatim}
UNTIL control until ( flag -- )

syntax: BEGIN <loop body> UNTIL
  effect:
    repeat <loop body> UNTIL
\end{verbatim}
\end{samepage}
These recipes identify exactly which control-word effects belong to the
one-versus-two calculation.  In particular, \texttt{WHILE} is included with
the loop prefix and \texttt{UNTIL} with the loop body.

\subsection{Numeric, logical, memory, and input words}
\begin{verbatim}
+ ( n n -- n )
- ( n n -- n )
0= ( x -- flag )
OR ( flag flag -- flag )
INVERT ( flag -- flag )

@ ( a-addr -- x )
C@ ( c-addr -- char )
KEY ( -- char )
\end{verbatim}
This profile lets \texttt{0=} accept \texttt{x}.  The presentation uses the
narrower input type \texttt{n}.  Section~\ref{sec:control} explains why both
choices give the same result here.

\begin{thebibliography}{99}
\footnotesize
\setlength{\itemsep}{0pt}
\setlength{\parsep}{0pt}
\raggedright

\bibitem{cousot1977}
P. Cousot and R. Cousot.
\newblock Abstract interpretation: A unified lattice model for static analysis
of programs by construction or approximation of fixpoints.
\newblock In \emph{Proceedings of POPL}, pp. 238--252, 1977.
\newblock doi:10.1145/512950.512973.
\newblock Public full text:
\url{https://www.di.ens.fr/~cousot/publications.www/CousotCousot-POPL-77-ACM-p238--252-1977.pdf}.

\bibitem{forth2012}
Forth 200x Standardisation Committee.
\newblock \emph{Forth 2012 Standard}.
\newblock Draft proposed standard, 10 November 2014.
\newblock Public full text:
\url{https://forth-standard.org/standard/intro}.

\bibitem{leroy2003}
X. Leroy.
\newblock Java bytecode verification: Algorithms and formalizations.
\newblock \emph{Journal of Automated Reasoning}, 30:235--269, 2003.
\newblock doi:10.1023/A:1025055424017.
\newblock Public full text:
\url{https://xavierleroy.org/publi/bytecode-verification-JAR.pdf}.

\bibitem{poial1994}
J. Pöial.
\newblock Algebraic specification of stack effects.
\newblock \emph{Forth Dimensions}, 16(4):18--20, 1994.
\newblock Public full text:
\url{https://www.forth.org/fd/FD-V16N4.pdf}.

\bibitem{poial1992}
J. Pöial.
\newblock Multiple stack-effects of Forth programs.
\newblock In \emph{1991 FORML Conference Proceedings, euroFORML'91},
pp. 400--406, 1992.
\newblock Public full text:
\url{https://www.complang.tuwien.ac.at/anton/euroforth/ef91/poial.pdf}.

\bibitem{poial2002}
J. Pöial.
\newblock Stack effect calculus with typed wildcards, polymorphism and
inheritance.
\newblock Presentation and proceedings abstract, 18th EuroForth Conference,
2002.
\newblock Public full text:
\url{https://www.complang.tuwien.ac.at/anton/euroforth/ef02/poial02.pdf}.

\bibitem{poial2006}
J. Pöial.
\newblock Typing tools for typeless stack languages.
\newblock In \emph{Proceedings of the 22nd EuroForth Conference},
pp. 40--46, 2006.
\newblock Public full text:
\url{https://www.complang.tuwien.ac.at/anton/euroforth/ef06/poial06.pdf}.

\bibitem{poial2008}
J. Pöial.
\newblock Java framework for static analysis of Forth programs.
\newblock In \emph{Proceedings of the 24th EuroForth Conference},
pp. 20--23, 2008.
\newblock Public full text:
\url{https://www.complang.tuwien.ac.at/anton/euroforth/ef08/papers/poial.pdf}.

\bibitem{saabas2006}
A. Saabas and T. Uustalu.
\newblock Compositional type systems for stack-based low-level languages.
\newblock In \emph{Proceedings of CATS 2006}, CRPIT 51, pp. 27--39, 2006.
\newblock Public full text:
\url{https://web.archive.org/web/20120325100058id_/http://crpit.com/confpapers/CRPITV51Saabas.pdf}.

\bibitem{stoddart1991}
B. Stoddart and P. J. Knaggs.
\newblock Type inference in stack based languages.
\newblock In \emph{1991 FORML Conference Proceedings, euroFORML'91},
pp. 407--421, 1991.
\newblock Public full text:
\url{https://www.complang.tuwien.ac.at/anton/euroforth/ef91/stoddart.pdf}.

\end{thebibliography}

\end{document}
